% Created 2026-08-18 ter 16:58
% Intended LaTeX compiler: pdflatex
\documentclass[11pt]{article}
\usepackage[utf8]{inputenc}
\usepackage[T1]{fontenc}
\usepackage{graphicx}
\usepackage{longtable}
\usepackage{wrapfig}
\usepackage{rotating}
\usepackage[normalem]{ulem}
\usepackage{amsmath}
\usepackage{amssymb}
\usepackage{capt-of}
\usepackage{hyperref}
\author{censo-escolar--analysis}
\date{\today}
\title{Plano: mapas dos microdados em HTML estático (GitHub Pages)}
\hypersetup{
 pdfauthor={censo-escolar--analysis},
 pdftitle={Plano: mapas dos microdados em HTML estático (GitHub Pages)},
 pdfkeywords={},
 pdfsubject={},
 pdfcreator={Emacs 30.2 (Org mode 9.7.11)}, 
 pdflang={Pt-Br}}
\begin{document}

\maketitle
\setcounter{tocdepth}{2}
\tableofcontents

Este documento é o plano dos mapas. Ele foi consolidado em 2026-08-18 depois
de três rodadas de conversa; o \emph{Histórico das decisões}, no fim, guarda o
caminho percorrido — inclusive as ideias descartadas e o motivo.

A ideia central cabe em uma linha de comando:

\begin{verbatim}
make map QT_MAT_BAS "https://servicodados.ibge.gov.br/api/v3/malhas/estados/35?formato=application/vnd.geo+json&qualidade=minima&intrarregiao=microrregiao" 2023
\end{verbatim}

\textbf{Uma variável do INEP + uma URL de malha do IBGE = um mapa publicado.} Todo
o resto deste documento explica as regras que transformam esses dois
argumentos numa página HTML.
\section{Objetivo}
\label{sec:org23efada}

Publicar mapas dos microdados do Censo Escolar como HTMLs estáticos no
GitHub Pages:

\begin{itemize}
\item um \texttt{index.html} que lista os mapas disponíveis;
\item cada mapa é uma página que carrega um único arquivo de dados e desenha,
sem servidor por trás;
\item os mapas são \emph{assados} por comandos \texttt{make} reproduzíveis, e a mescla de
vários anos usa a coluna \texttt{NU\_ANO\_CENSO};
\item o visitante não monta consultas: ele escolhe um mapa da lista. Quem monta
a consulta é quem roda o \texttt{make}.
\end{itemize}
\section{O comando}
\label{sec:org370c2e2}

\begin{verbatim}
make map <VARIAVEL> "<URL da malha>" [anos...] [OPCAO=valor...]
\end{verbatim}

\begin{center}
\begin{tabular}{lll}
Argumento & Obrigatório & Exemplo\\
\hline
variável do INEP & sim & \texttt{QT\_MAT\_BAS}, \texttt{IN\_INTERNET}\\
URL da malha IBGE & sim & \texttt{"https://…/estados/35?…\&intrarregiao=microrregiao"}\\
ano(s) & não & \texttt{2023}, ou \texttt{2019 2024} (padrão: o mais recente que serve)\\
\texttt{CATEGORIA} & só p/ \texttt{TP\_} & \texttt{CATEGORIA=4} (privada)\\
\texttt{TITULO} & não & \texttt{TITULO}"Matrículas na básica — SP"=\\
\end{tabular}
\end{center}

\textbf{As aspas na URL são obrigatórias.} Isso é do shell, não do \texttt{make}: sem
elas, o \texttt{\&} da URL manda o comando para segundo plano e nada roda.

Exemplos:

\begin{verbatim}
make map QT_MAT_BAS  "https://…/estados/35?…&intrarregiao=microrregiao" 2023
make map IN_INTERNET "https://…/paises/BR?…&intrarregiao=UF" 2019 2024
make map TP_DEPENDENCIA "https://…/estados/31?…&intrarregiao=municipio" 2023 CATEGORIA=4
\end{verbatim}
\section{As regras de geração}
\label{sec:org95ec85a}

Esta é a seção central: como os dois argumentos viram um mapa. Cada regra
tem a evidência que a sustenta — tudo abaixo foi verificado na API do IBGE e
nos microdados em disco, em 2026-08-18.
\subsection{R1 — A URL carrega três informações, e o gerador lê as três}
\label{sec:orgf641db7}

\begin{verbatim}
https://servicodados.ibge.gov.br/api/v3/malhas/estados/35?formato=…&qualidade=minima&intrarregiao=microrregiao
                                             └─────┬────┘            └───┬────┘        └────────┬─────────┘
                                          recorte espacial         custo/detalhe        grão do mapa
\end{verbatim}

\begin{center}
\begin{tabular}{ll}
Parte da URL & Vira\\
\hline
caminho (\texttt{/estados/35}) & \textbf{filtro} das linhas dos microdados\\
\texttt{intrarregiao} & \textbf{coluna de agrupamento} da agregação\\
\texttt{qualidade} & tamanho do arquivo (só afeta o peso)\\
\texttt{formato} & normalizado pelo gerador (veja R9)\\
\end{tabular}
\end{center}

O caminho vira filtro assim:

\begin{center}
\begin{tabular}{ll}
Caminho & Filtro aplicado\\
\hline
\texttt{/paises/BR} & nenhum (Brasil inteiro)\\
\texttt{/estados/35} & \texttt{CO\_UF =} 35=\\
\texttt{/municipios/3550308} & \texttt{CO\_MUNICIPIO =} 3550308=\\
\end{tabular}
\end{center}
\subsection{R2 — \texttt{intrarregiao} determina a coluna de agrupamento}
\label{sec:org1a564c3}

Tabela verificada na API (contagem de feições e formato de \texttt{codarea} reais):

\begin{center}
\begin{tabular}{lllrl}
\texttt{intrarregiao} & \texttt{codarea} & Coluna agrupadora & Feições & Anos em que a coluna existe\\
\hline
(ausente, em \texttt{/estados/35}) & \texttt{35} & \texttt{CO\_UF} & 1 & todos\\
\texttt{UF} & \texttt{11} & \texttt{CO\_UF} & 27 & todos\\
\texttt{regiao-intermediaria} & \texttt{1102} & \texttt{CO\_REGIAO\_GEOG\_INTERM} & 133 & 2023–2025\\
\texttt{mesorregiao} & \texttt{1102} & \texttt{CO\_MESORREGIAO} & 137 & 2019–2024\\
\texttt{regiao-imediata} & \texttt{110005} & \texttt{CO\_REGIAO\_GEOG\_IMED} & 510 & 2023–2025\\
\texttt{microrregiao} & \texttt{35035} & \texttt{CO\_MICRORREGIAO} & 558 & 2019–2024\\
\texttt{municipio} & \texttt{1100015} & \texttt{CO\_MUNICIPIO} & 5.570 & todos\\
\end{tabular}
\end{center}

\textbf{Armadilha registrada:} mesorregião e região intermediária usam ambas
códigos de 4 dígitos — \texttt{1102} é uma de cada, e são divisões diferentes. Não
dá para inferir o grão pelo formato do código; tem de se ler o
\texttt{intrarregiao} da URL. É um bom motivo a mais para a URL ser a consulta.

\textbf{Limite da API} (testado, o erro é explícito): em \texttt{/estados/XX} só valem
\texttt{mesorregiao}, \texttt{microrregiao} e \texttt{municipio}; \texttt{/municipios/XXXXXXX} não
aceita \texttt{intrarregiao} nenhum. Distrito não é servido por esta API — veja
"Até onde o drill-down chega".
\subsection{R3 — A junção é direta: \texttt{codarea} ↔ coluna agrupadora}
\label{sec:org33f07c8}

O \texttt{properties.codarea} de cada polígono casa com o valor da coluna
agrupadora \textbf{sem conversão}. Conferido na URL de exemplo: a malha devolve
\texttt{codarea = "35035"} e uma escola de São Paulo em 2023 traz
\texttt{CO\_MICRORREGIAO = 35061} — mesmo formato, mesma origem (IBGE).

O gerador compara os dois conjuntos e \textbf{relata a cobertura}: quantos
polígonos ficaram sem dado e quantos códigos de dado não acharam polígono.
Um mapa com buracos silenciosos é pior que um erro.
\subsection{R4 — A variável decide a agregação e o tipo de mapa}
\label{sec:org98b4cd0}

Quem responde isso é o dicionário de dados que o projeto \textbf{já gera}
(\texttt{censo dicionario} → \texttt{data/processed/dicionario\_colunas\_comuns.csv}), na
coluna \texttt{natureza}. Contagem do arquivo atual:

\begin{center}
\begin{tabular}{lrlll}
\texttt{natureza} & Nº & Exemplo & Agregação por área & Mapa\\
\hline
contagem (\texttt{QT\_}) & 168 & \texttt{QT\_MAT\_BAS} & soma & coroplético do total\\
indicador 0/1 (\texttt{IN\_}) & 282 & \texttt{IN\_INTERNET} & média → \% das escolas & coroplético de taxa (0–100 \%)\\
categórica (\texttt{TP\_}) & 17 & \texttt{TP\_DEPENDENCIA} & \% da categoria pedida & coroplético da categoria\\
código/nome/data/etc. & 34 & \texttt{CO\_UF}, \texttt{NO\_*} & — & \textbf{recusa, com explicação}\\
\end{tabular}
\end{center}

São \textasciitilde{}467 variáveis que viram mapa sem código novo. Duas ressalvas que o
gerador aplica sozinho:

\begin{itemize}
\item \texttt{TP\_*} exige \texttt{CATEGORIA=}; sem ela, o padrão é pintar a \textbf{categoria
dominante} de cada área;
\item toda página ganha, no rodapé, o \texttt{descricao} oficial do INEP para aquela
variável — texto que o dicionário já traz.
\end{itemize}
\subsection{R5 — Os anos filtram \texttt{NU\_ANO\_CENSO}, e a disponibilidade é conferida antes}
\label{sec:org50e77ba}

Anos soltos na linha de comando (\texttt{2023}, ou \texttt{2019 2024}) viram filtro sobre
\texttt{NU\_ANO\_CENSO} depois do empilhamento feito por \texttt{loading.carregar\_anos()}.
Sem ano, usa-se o mais recente em que a variável \textbf{e} a coluna agrupadora
existem.

Antes de ler qualquer CSV, o gerador cruza três disponibilidades e recusa
com mensagem clara se alguma falhar:

\begin{enumerate}
\item a \textbf{variável} existe naquele ano? (o dicionário sabe: \texttt{anos\_ausentes});
\item a \textbf{coluna agrupadora} existe naquele ano? (tabela da R2 — pedir
\texttt{microrregiao} em 2025 ou \texttt{regiao-imediata} em 2019 é erro);
\item o \textbf{tipo é estável} entre os anos pedidos? Se
\texttt{tipo\_estavel = False}, o mapa é gerado com uma ressalva no rodapé.
\end{enumerate}

Com mais de um ano, a página ganha um seletor e o arquivo de dados traz uma
série por área — nada de baixar de novo ao trocar o ano.
\subsection{R6 — O nome do mapa é determinístico}
\label{sec:orgba1002e}

\texttt{<variavel>-<intrarregiao>-<recorte>-<anos>}, em minúsculas. A URL de
exemplo com \texttt{QT\_MAT\_BAS} e 2023 produz:

\begin{verbatim}
site/mapas/qt_mat_bas-microrregiao-35-2023.html
site/dados/qt_mat_bas-microrregiao-35-2023.geojson
\end{verbatim}

Rodar o mesmo comando de novo \textbf{sobrescreve} a mesma página, em vez de
acumular lixo. É o que permite reassar o site inteiro sem limpar nada antes.
\subsection{R7 — O produto é um GeoJSON com o dado já dentro}
\label{sec:org61c9b3c}

Aqui está a simplificação que mais economiza trabalho, e ela decorre de uma
decisão sua: \emph{as malhas baixadas ficam no \texttt{.gitignore}}. Se a malha crua não
entra no git, ela não pode ser servida ao navegador — então o gerador
\textbf{funde} dado e geometria numa coisa só:

\begin{verbatim}
{"type": "FeatureCollection",
 "features": [
   {"type": "Feature",
    "properties": {"codarea": "35035", "nome": "Adamantina",
                   "valor": {"2023": 148}, "escolas": {"2023": 148}},
    "geometry": {...}}
 ],
 "meta": {"variavel": "QT_MAT_BAS", "descricao": "Número de Matrículas…",
          "natureza": "contagem", "anos": [2023], "fonte_malha": "<url>",
          "gerado_em": "2026-08-18", "cobertura": "63/63 áreas com dado"}}
\end{verbatim}

Consequências:

\begin{itemize}
\item \textbf{uma requisição por página} (o GeoJSON), em vez de duas (malha + dados);
\item só entram no git as áreas efetivamente usadas — o Brasil inteiro por
município nunca é servido quando o mapa é de uma UF;
\item o \texttt{meta} vai junto, então a página se documenta sozinha (fonte, data de
geração, cobertura da junção) sem arquivo paralelo;
\item o JavaScript encolhe para "buscar um arquivo, pintar, mostrar popup".
\end{itemize}

E, do mesmo comando, sai também o \textbf{CSV agregado} em
\texttt{data/processed/<slug>.csv} — a tabela que gerou o mapa, conferível no Calc
ou com \texttt{censo amostra -{}-{}arquivo}. O CSV existe para inspeção e reuso; o
GeoJSON, para o navegador. Saem juntos justamente para nunca divergirem: se
o número do mapa parecer errado, a tabela que o produziu está ali do lado.
\subsection{R8 — O \texttt{make map} também registra a consulta}
\label{sec:orgfac8044}

Cada geração acrescenta (ou atualiza) uma linha em \texttt{site/consultas.txt}, com
exatamente o que se digitou:

\begin{verbatim}
# variável     URL da malha                                          anos   opções
QT_MAT_BAS     "https://…/estados/35?…&intrarregiao=microrregiao"    2023
IN_INTERNET    "https://…/paises/BR?…&intrarregiao=UF"               2019 2024
\end{verbatim}

\textbf{O site inteiro é função desse arquivo.} \texttt{make site} relê a lista e reassa
tudo; o \texttt{index.html} é gerado dela, nunca editado à mão; apagar um mapa é
apagar uma linha; e o \texttt{git diff} mostra literalmente quais mapas entraram,
saíram ou mudaram. Um mapa publicado passa a ser revisável como código.
\subsection{R9 — As regras de robustez}
\label{sec:org4a1439f}

\begin{enumerate}
\item \textbf{Normalizar o \texttt{formato=}}: se a URL vier sem
\texttt{formato=application/vnd.geo+json}, o gerador acrescenta — sem isso a API
devolve outra representação e o erro só apareceria ao desenhar.
\item \textbf{Cache por hash da URL} em \texttt{data/cache/malhas/<sha1>.json} (no
\texttt{.gitignore}): reassar o site não depende da API no ar, e não se martela
o IBGE a cada build.
\item \textbf{Avisar quando a malha for pesada}: acima de \textasciitilde{}150 KiB gzipados, sugerir
\texttt{qualidade=minima}. Medido nas mesmas 63 microrregiões de SP:

\begin{center}
\begin{tabular}{lll}
\texttt{qualidade} & Bruto & Gzip\\
\hline
\texttt{maxima} & 1,77 MiB & \textbf{416 KiB}\\
\texttt{minima} & 91 KiB & \textbf{21,7 KiB}\\
\end{tabular}
\end{center}

Dezenove vezes menor, pelos mesmos polígonos. A malha é fundo de mapa: a
precisão de vértice não muda a leitura do dado. É aviso, não imposição —
num mapa de uma cidade só, \texttt{maxima} pode fazer sentido.
\item \textbf{Detectar a URL sem aspas}: se chegar uma URL truncada (sem
\texttt{intrarregiao=/=qualidade}), dizer "põe a URL entre aspas" em vez de
falhar torto.
\item \textbf{Falhar antes de processar}: toda validação acontece na leitura do
Makefile ou nas primeiras linhas do comando — nunca depois de varrer
200 MiB. É a mesma disciplina que já aplicamos em \texttt{make dicionario}.
\end{enumerate}
\subsection{Exemplo completo, do comando à página}
\label{sec:org8c48921}

Para \texttt{make map QT\_MAT\_BAS "https://…/estados/35?…\&intrarregiao=microrregiao" 2023}:

\begin{center}
\begin{tabular}{rll}
\# & O que acontece & Regra\\
\hline
1 & \texttt{\$(filter http\%,\$(MAKECMDGOALS))} captura a URL; \texttt{2023} vira o ano & R1/R5\\
2 & \texttt{QT\_MAT\_BAS} → natureza \texttt{contagem} → agregação = soma & R4\\
3 & \texttt{intrarregiao=microrregiao} → agrupar por \texttt{CO\_MICRORREGIAO} (existe em 2023) & R2/R5\\
4 & \texttt{/estados/35} → filtrar \texttt{CO\_UF =} 35= & R1\\
5 & \texttt{carregar\_anos([2023], colunas}{[}NU\textsubscript{ANO}\textsubscript{CENSO}, CO\textsubscript{UF}, CO\textsubscript{MICRORREGIAO}, QT\textsubscript{MAT}\textsubscript{BAS}])= & —\\
6 & \texttt{groupby(CO\_MICRORREGIAO).sum()} → 63 linhas & R4\\
7 & malha baixada (cache) e unida por \texttt{codarea}; cobertura relatada & R3/R9\\
8 & grava \texttt{site/dados/qt\_mat\_bas-microrregiao-35-2023.geojson} + \texttt{.html} & R6/R7\\
9 & acrescenta a linha em \texttt{site/consultas.txt} e regenera o \texttt{index.html} & R8\\
\end{tabular}
\end{center}
\section{Arquitetura e arquivos}
\label{sec:org819c827}

\begin{verbatim}
microdados (data/interim, data/processed)          [já existe]
      │
      │  make map <VAR> "<URL>" [anos]             [Python: censo map]
      │      ├─ malha do IBGE  → data/cache/malhas/<sha1>.json   (.gitignore)
      │      └─ agregação      → junção por codarea
      ▼
site/dados/<slug>.geojson     ← geometria + dado + meta, num arquivo só  (git)
site/mapas/<slug>.html        ← página gerada do modelo                  (git)
site/consultas.txt            ← a lista que define o site                (git)
site/index.html               ← gerado da lista                          (git)
      │
      │  GitHub Actions publica site/
      ▼
https://wagnermarques.github.io/censo-escolar--analysis/
\end{verbatim}
\subsection{O que entra no git e o que não entra}
\label{sec:org2a572c6}

\begin{center}
\begin{tabular}{lll}
Caminho & Git & Por quê\\
\hline
\texttt{data/cache/malhas/} & não & download bruto, grande, reproduzível\\
\texttt{data/raw}, \texttt{interim} & não & microdados, já ignorados hoje\\
\texttt{data/processed/*.csv} & não & tabela agregada, inspecionável, reproduzível\\
\texttt{site/dados/*.geojson} & sim & é o que o Pages serve; pequeno e recortado\\
\texttt{site/mapas/*.html} & sim & gerado, mas é o produto publicado\\
\texttt{site/consultas.txt} & sim & \textbf{o código-fonte do site}\\
\end{tabular}
\end{center}
\subsection{JavaScript}
\label{sec:orga7fa1c7}

Só o essencial, vendorado em \texttt{site/vendor/} (sem CDN, sem build):

\begin{center}
\begin{tabular}{lll}
Biblioteca & Papel & Por quê\\
\hline
\textbf{Leaflet} & tiles, GeoJSON, popup, zoom & \textasciitilde{}42 KiB gz; API estável há uma década\\
\texttt{fetch()} nativo & carregar o \texttt{.geojson} & zero dependência\\
JS vanilla & escala de cor, legenda, seletor de ano & \textasciitilde{}100 linhas, num modelo único p/ todos os mapas\\
\end{tabular}
\end{center}

Ficam de fora, e o motivo: \emph{PapaParse} (o Python já entrega o arquivo
pronto), \emph{D3} (o Leaflet resolve coroplético com um décimo do que há para
aprender) e qualquer \emph{framework com build} (a exigência é \texttt{git push} e
pronto). Única dependência externa em runtime: os tiles do OpenStreetMap.
\section{Comandos make}
\label{sec:orgc0e74e1}

\begin{verbatim}
make map <VAR> "<URL>" [anos] [OPCAO=v]   # assa um mapa e registra a consulta
make site                                 # reassa tudo a partir de site/consultas.txt
make servir                               # http.server em site/, para ver localmente
make ibge-malhas                          # pré-aquece o cache com as malhas mais usadas
make ibge-malha-municipal                 # idem, só a malha municipal (a mais pesada)
\end{verbatim}

\texttt{make ibge-malhas} é conveniência, não obrigação: o \texttt{make map} baixa o que
faltar. Ele existe para quem vai trabalhar sem rede depois, e as malhas
podem ser consultadas e baixadas à mão em
\url{https://www.ibge.gov.br/geociencias/organizacao-do-territorio/malhas-territoriais/15774-malhas.html}
(Municípios, UFs, Grandes Regiões, Regiões Geográficas Imediatas e
Intermediárias, País). O que ele baixa vai para \texttt{data/cache/malhas/}, que
está no \texttt{.gitignore}.

Cada alvo com sua entrada em \texttt{make help <alvo>}, e por baixo os subcomandos
\texttt{censo map}, \texttt{censo site} e \texttt{censo servir} — a paridade make/CLI que o
projeto já pratica (funciona no Windows sem \texttt{make}).
\section{Fatos verificados}
\label{sec:org609f35e}

Tudo conferido em 2026-08-18, nos dados em disco e nas fontes oficiais.
\subsection{Os microdados não têm latitude/longitude}
\label{sec:org7c2209f}

Não existe coluna \texttt{LATITUDE=/=LONGITUDE} em nenhum ano. O que há é
\texttt{CO\_MUNICIPIO}, UF, endereço em texto e CEP. Coordenadas por escola existem
no \href{https://www.gov.br/inep/pt-br/acesso-a-informacao/dados-abertos/inep-data/catalogo-de-escolas}{Catálogo de Escolas}, cujo export é um portal de BI \textbf{sem URL estável de
download} — não dá para automatizar no \texttt{make}, e nem toda escola tem
coordenada lá.

\emph{Decisão registrada:} mapas por área (UF, região, município), sem pontos por
escola. É o que os dados permitem hoje.
\subsection{Até onde o drill-down chega}
\label{sec:orgb6ccc80}

Você pediu para descer de estado → região → distrito. O que dá e o que não
dá:

\textbf{Nos dados} — colunas por ano:

\begin{center}
\begin{tabular}{ll}
Coluna & Anos\\
\hline
\texttt{CO\_REGIAO}, \texttt{CO\_UF}, \texttt{CO\_MUNICIPIO} & todos (2019–2025)\\
\texttt{CO\_MESORREGIAO}, \texttt{CO\_MICRORREGIAO} & 2019–2024 (some em 2025)\\
\texttt{CO\_REGIAO\_GEOG\_INTERM} / \texttt{\_IMED} & 2023–2025\\
\texttt{CO\_DISTRITO} & \textbf{todos} (2019–2025)\\
\texttt{CO\_SUBDISTRITO} & só 2025\\
\end{tabular}
\end{center}

\texttt{CO\_DISTRITO} é o código IBGE de 9 dígitos (município + 2): \texttt{355030840} em
São Paulo. \textbf{O dado desce até o distrito em todos os anos.}

\textbf{Na geometria}: a API de malhas \textbf{não serve distrito} (testado). Existe outra
fonte — \texttt{geoftp.ibge.gov.br/…/malhas\_de\_setores\_censitarios\_\_divisoes\_intramunicipais/}
—, em shapefile, o que exigiria uma dependência nova (geopandas/fiona) e um
alvo \texttt{make} à parte.

\textbf{Os nomes} dos distritos, esses são fáceis:
\texttt{api/v1/localidades/estados/35/distritos} devolve os 1.042 distritos de SP
com id de 9 dígitos e nome — junta direto com \texttt{CO\_DISTRITO}, sem
dependência nova.

\emph{Portanto:} distrito entra primeiro como \textbf{painel de dados} (ranking, tabela,
barras dentro de um município), não como polígono. O polígono fica para uma
fase posterior, se valer a pena.

Sobre "zona leste", "zona norte", "centro": não são divisões do IBGE — não
existem nos microdados nem nas malhas. As oficiais em São Paulo são as
\textbf{subprefeituras}, agrupamentos de distritos definidos pela prefeitura.
Podemos ter isso como um arquivo de agrupamento nosso
(\texttt{site/dados/agrupamentos/sp-zonas.json}: distrito → zona), assumido como
convenção do projeto, não como dado oficial.
\subsection{O buraco de 2025}
\label{sec:orgfee20d9}

A \texttt{Tabela\_Escola\_2025\_V2.csv} \textbf{não tem nenhuma coluna} \texttt{QT\_MAT\_*}: em 2025 o
INEP moveu as matrículas para a \texttt{Tabela\_Matricula}, um arquivo à parte.
Então: contagem de escolas cobre 2019–2025; matrículas, 2019–2024. O
gerador recusa \texttt{QT\_MAT\_BAS} em 2025 com essa explicação (R5).
\subsection{Custo das malhas (medido, gzip que o Pages já serve)}
\label{sec:org8dbdea6}

\begin{center}
\begin{tabular}{lll}
Malha & Bruto & Gzip\\
\hline
BR por UF & 98 KiB & 28 KiB\\
BR por região intermediária & 613 KiB & 163 KiB\\
BR por região imediata & 1,1 MiB & 283 KiB\\
BR por município (5.570) & 3,6 MiB & 817 KiB\\
SP por município & 310 KiB & 67 KiB\\
SP por microrregião (\texttt{minima}) & 91 KiB & 21,7 KiB\\
SP por microrregião (\texttt{maxima}) & 1,77 MiB & 416 KiB\\
\end{tabular}
\end{center}

A leitura prática: \textbf{recortar por UF é o que barateia tudo} (67 KiB contra
817 KiB), e \texttt{qualidade=minima} é o padrão sensato.
\subsection{O que o projeto já tem e será reaproveitado}
\label{sec:org29afb27}

\begin{center}
\begin{tabular}{ll}
Já existe & Serve para\\
\hline
\texttt{loading.carregar\_anos()} & empilhar anos com \texttt{NU\_ANO\_CENSO} (R5)\\
\texttt{loading.carregar\_escolas()} & ler só as colunas necessárias (Parquet/CSV)\\
\texttt{analysis.contar\_escolas()} & a agregação por área (R4)\\
\texttt{esquema} / dicionário de dados & natureza, descrição, anos ausentes (R4/R5)\\
\texttt{codigos.rotular()} & rótulos de \texttt{TP\_*} nos popups\\
\texttt{amostra.py} & conferir o CSV intermediário sem abrir tudo\\
modificadores do Makefile & anos soltos e validação antecipada (R1/R5/R9)\\
padrão \texttt{make help <alvo>} & ajuda detalhada do \texttt{make map} de graça\\
\end{tabular}
\end{center}
\section{Decisões registradas}
\label{sec:orgad68be9}

\begin{enumerate}
\item \textbf{Deploy}: \texttt{site/} na \texttt{main} + GitHub Actions (workflow de \textasciitilde{}15 linhas com
\texttt{upload-pages-artifact} + \texttt{deploy-pages}). \emph{Escolhido pelo usuário.}
\item \textbf{Sem pontos por escola} no v1 — mapas por área, pelos dados que temos.
\item \textbf{Malhas baixadas ficam no \texttt{.gitignore}}; o que é publicado é o GeoJSON
recortado com o dado dentro (R7).
\item \textbf{O construtor de consultas é o \texttt{make}}, não uma tela. O visitante escolhe
um mapa pronto na lista.
\item \textbf{A URL do IBGE é a linguagem do recorte} — o projeto não inventa
vocabulário próprio de nível geográfico.
\end{enumerate}
\section{Fases}
\label{sec:org48d4936}

\begin{itemize}
\item \textbf{F1} — \texttt{make map QT\_MAT\_BAS "…/paises/BR?…\&intrarregiao=UF" 2023} assa a
primeira página, o \texttt{index.html} e o workflow do Pages. \emph{Aqui o objetivo já
está no ar.}
\item \textbf{F2} — \texttt{site/consultas.txt} + \texttt{make site} reassando tudo; mais mapas.
\item \textbf{F3} — dentro da página: seletor de ano (quando a consulta tem vários) e
filtro por dependência, com o dado que o GeoJSON já traz.
\item \textbf{F4} — distrito como painel de dados (nomes via API de localidades).
\item \textbf{F5} — se valer: polígono de distrito (shapefile do geoftp) e/ou pontos
por escola (Catálogo, export manual).
\end{itemize}
\section{Riscos e limites}
\label{sec:org2e8baf7}

\begin{itemize}
\item \textbf{Dependência da API do IBGE} no momento da geração — mitigada pelo cache
(R9.2), que torna \texttt{make site} reprodutível offline.
\item \textbf{Tiles do OSM} são a única dependência externa em runtime; trocar de
provedor é uma linha.
\item \textbf{Malha pesada} se alguém pedir Brasil inteiro por município em
\texttt{qualidade=maxima} — o aviso da R9.3 existe para isso.
\item \textbf{Junção incompleta} entre malha e dados (municípios criados/extintos entre
edições): por isso a cobertura é relatada, não presumida (R3).
\item \textbf{2025 sem matrículas} na tabela de escolas.
\end{itemize}
\section{Critérios de verificação}
\label{sec:org516ebe1}

\begin{itemize}
\item \texttt{make map QT\_MAT\_BAS "…\&intrarregiao=UF" 2023} gera os dois arquivos e o
total do GeoJSON bate com \texttt{analysis.panorama\_por\_uf()} para o mesmo ano;
\item a cobertura relatada é 27/27 áreas (UF) e 63/63 (microrregiões de SP);
\item \texttt{make servir} + navegador: o índice lista o mapa, o mapa desenha, o popup
mostra nome e valor;
\item \texttt{make site} a partir de um \texttt{consultas.txt} com três linhas reproduz três
páginas idênticas às geradas uma a uma;
\item depois do push, a URL do Pages mostra o mesmo resultado.
\end{itemize}
\section{Histórico das decisões}
\label{sec:org0bad18a}

Três rodadas de conversa, e o que cada uma descartou:

\begin{enumerate}
\item \textbf{Páginas fixas} (proposta inicial): um HTML por tema, escrito à mão.
Descartado por não escalar para as \textasciitilde{}467 variáveis mapeáveis.
\item \textbf{Construtor de consultas na tela} (ideia do usuário, 2026-08-18): a URL
da página como consulta, com formulário gerado de um \texttt{manifesto.json}.
Descartado depois do esclarecimento seguinte — mas deixou duas ideias que
sobreviveram: a consulta como texto e o índice como lista de consultas
salvas.
\item \textbf{O construtor é o \texttt{make}} (esclarecimento do usuário): quem monta a
consulta é quem gera o site. Some o manifesto, some a UI, some o risco de
URL sem arquivo, some a explosão combinatória.
\item \textbf{A URL da malha é a consulta espacial} (proposta do usuário): a linguagem
do IBGE substitui a nossa. Some o vocabulário próprio de nível, some o
alvo \texttt{make malha} separado, e o comando fica com dois argumentos.
\end{enumerate}

O que sobreviveu de ponta a ponta: \emph{o site é função de uma lista de
consultas}, e \emph{o Python calcula, o navegador só desenha}.
\end{document}
